\documentclass[11pt,a4paper]{article}

% ── Packages ──────────────────────────────────────────────────────────────
\usepackage[utf8]{inputenc}
\usepackage[T1]{fontenc}
\usepackage{amsmath,amssymb,amsthm}
\usepackage{amscd}
\usepackage{booktabs}
\usepackage{graphicx}
\usepackage{hyperref}
\usepackage[margin=1in]{geometry}
\usepackage{enumitem}
\usepackage{xcolor}
\definecolor{hlblue}{RGB}{0,80,160}
\hypersetup{colorlinks=true,linkcolor=hlblue,urlcolor=hlblue,citecolor=hlblue}

% ── Theorems ──────────────────────────────────────────────────────────────
\newtheorem{principle}{Principle}[section]
\newtheorem{theorem}{Theorem}[section]
\newtheorem{definition}{Definition}[section]
\newtheorem{proposition}{Proposition}[section]
\theoremstyle{remark}
\newtheorem{remark}{Remark}[section]

% ── Title ─────────────────────────────────────────────────────────────────
\title{\textbf{Emerging World Models}\\[4pt]
\large A Structural Framework for Self-Organizing World Representations
\\[4pt]
\large \textsc{Consolidation Document --- August 2026}}
\author{Alex Mylnikov\\[4pt]
\texttt{email: alexmy@lisa-park.com}}
\date{}

\begin{document}
\maketitle

\begin{abstract}
We propose \textbf{Emerging World Models} (EWM)---a structural framework for world
representations whose ontology is not predefined but crystallizes through measurement.
EWM is governed by four foundation principles. \textbf{(1) Emergent Ontology:} elements
and subsystems are not primitive; they emerge from the accumulation of IICA measurement
interactions. \textbf{(2) IICA Morphisms:} every relation between elements is an
Idempotent, Immutable, Content-Addressed morphism; composition preserves these
properties, enabling arbitrarily nested world representations without new theory.
\textbf{(3) The Ashby--Bootstrap Correspondence:} Ashby's Requisite Variety Law,
interpreted naively, demands cardinality balance between connected systems. EWM resolves
this not by matching cardinalities directly, but through \textit{bootstrap
permutation}---multiplying measurements (e.g., n-grams) to bridge representational gaps
without sacrificing resolution. \textbf{(4) Noether Evolution:} structural change in an
EWM obeys conservation laws; every symmetry of the measurement regime yields a conserved
quantity. The lattice union invariant, the DRN decomposition, and the holographic
reconstruction property are all consequences of this principle. We ground each principle
in the concrete instantiation of HLLSet Algebra and articulate the target formalism---a
Grothendieck topology on the category of measurement patches---toward which the framework
aspires, in spirit if not yet in letter.
\end{abstract}

% ==========================================================================
\section{Introduction}
% ==========================================================================

A world model is a system that learns an internal representation of its environment
from sensory input, then plans and reasons within that representation. Classical
approaches---JEPA, large language models, and token-based architectures---share a common
presupposition: the elements of the world (tokens, objects, features) are \textit{given}
before the model begins to operate. The vocabulary is fixed. The ontology is predefined.
The model's task is to learn the relations among these predefined elements.

\textbf{Emerging World Models} (EWM) invert this presupposition. In an EWM, elements are
not predefined. They \textit{emerge} from the accumulation of measurement interactions.
An element is not a token in a vocabulary; it is a stable pattern that crystallizes when
repeated measurements converge. Subsystems are not architectural modules; they are
coherent clusters of elements that co-occur across measurement trajectories.

This inversion has deep structural consequences. It shifts the foundation of world
modeling from \textit{representation} (encoding predefined elements into a vector space)
to \textit{measurement} (accumulating observations into a holographic plate from which
elements can be recovered). It replaces mutable, named parameters with immutable,
content-addressed bitmasks. It transforms learning from gradient descent over weight
matrices to rank rearrangement in a content-addressed dictionary.

This document consolidates the theoretical principles discovered across the HLLSet
Algebra research program \cite{astesj2026,sgs2024} into a unified structural framework.
We articulate four principles that govern every EWM, ground them in the concrete
instantiation of HLLSet Algebra, and identify Grothendieck topology as the target
formalism toward which the framework aspires.

\subsection{Why ``Emerging''?}

The term ``emerging'' is deliberate. It signals three properties that distinguish EWM
from classical world models:

\begin{enumerate}[leftmargin=*,nosep]
  \item \textbf{Ontological emergence.} The identity of an element is not assigned by a
    designer. It is \textit{recognized} by the system when a measurement pattern
    stabilizes across repeated observations. A token that appears once is noise. A token
    that persists across temporal layers becomes an element.
  \item \textbf{Structural emergence.} Relations between elements are not wired by an
    architect. They are \textit{computed} as IICA morphisms between the content-addressed
    fingerprints of the elements. The lattice of relations emerges from the lattice of
    measurements.
  \item \textbf{Evolutionary emergence.} The system does not converge to a fixed point.
    It continuously reorganizes as new measurements arrive, old patterns fade, and the
    balance between novelty and retention shifts. Evolution is structural---governed by
    conservation laws---rather than parametric.
\end{enumerate}

% ==========================================================================
\section{Principle 1: Emergent Ontology}
\label{sec:emergence}
% ==========================================================================

\begin{principle}[Emergent Ontology]
\label{prin:emergence}
The elements of an Emerging World Model are not defined a priori. They emerge from
the accumulation of measurement interactions. An element is \textit{what persists}
across repeated observation.
\end{principle}

\subsection{The Classical Approach: Predefined Ontology}

In classical world models, the ontology is fixed before the model encounters any data:
\begin{itemize}[leftmargin=*,nosep]
  \item \textbf{LLMs} operate on a fixed token vocabulary (50K--250K tokens) defined by
    a tokenizer (BPE, SentencePiece). The tokenizer is trained once and frozen. Every
    subsequent observation is projected into this fixed vocabulary.
  \item \textbf{JEPA} operates on patch embeddings derived from a pretrained encoder.
    The patches are predefined spatial or temporal segments.
  \item \textbf{Symbolic AI} operates on hand-crafted symbol tables and ontologies.
\end{itemize}

In all cases, the system cannot \textit{discover} a new element. A concept that falls
between tokens is represented as a combination of existing tokens, never as a primitive
in its own right. A novel pattern that does not align with the pretrained encoder's
latent dimensions is invisible.

\subsection{The EWM Approach: Measurement as Ontogenesis}

In an EWM, the only primitive is the \textbf{measurement event}: a token entering the
system, hashed to a bit position, accumulating a mark in the holographic plate. There is
no ``vocabulary''---only a Lookup Table (LUT) that tracks which tokens have been observed
at which positions, and how often.

An element emerges when a measurement pattern satisfies two conditions:
\begin{enumerate}[leftmargin=*,nosep]
  \item \textbf{Stability.} The same token or token cluster maps to consistent bit
    positions across multiple measurement events.
  \item \textbf{Persistence.} The token's Term Frequency (TF) at its bit positions
    remains above a threshold across temporal layers.
\end{enumerate}

A token that appears once and never again leaves a mark (its bit is set in the HLLSet),
but it does not become an element---its TF is negligible, its rank is low, and it does
not participate in relations. A token that appears thousands of times across weeks
accumulates high TF, high rank, and dense R-links to co-occurring tokens. It has
\textit{emerged} as an element.

\subsection{Formalization}

Let $\mathcal{T}$ be the space of all possible tokens (e.g., all byte strings). Let
$h: \mathcal{T} \to \mathcal{P}$ be an IICA morphism mapping tokens to patches (bit
positions) in a finite patch space $\mathcal{P} = \{0,\ldots,N-1\}$.

\begin{definition}[Measurement Event]
A measurement event at time $t$ is a pair $(t_i, p_i)$ where $t_i \in \mathcal{T}$ is a
token and $p_i = h(t_i) \in \mathcal{P}$ is its patch assignment.
\end{definition}

\begin{definition}[Element]
An element $e \in \mathcal{E}$ is a token $t \in \mathcal{T}$ whose accumulated
measurement count $\text{TF}(t)$ exceeds an emergence threshold $\theta_e$ and whose
temporal distribution across layers is non-uniform (it has a ``lifetime,'' not just a
single appearance):
\begin{equation}
  e \in \mathcal{E} \iff \sum_{i=0}^{K} \text{TF}_i(t) > \theta_e
  \;\land\; \exists i,j: \text{TF}_i(t) \gg \text{TF}_j(t)
\end{equation}
\end{definition}

The set of elements $\mathcal{E}$ is not fixed. It grows as new tokens accumulate
measurements, and it contracts as tokens fade. The ontology is \textit{emergent} and
\textit{dynamic}.

\subsection{HLLSet Instantiation}

In HLLSet Algebra, the emergent ontology is realized through the LUT (Lookup Table).
The LUT is a monotonic CRDT that maps each (register, trailing\_zeros) bucket to the
set of tokens that hash to that bucket, together with their accumulated TF values. The
LUT \textit{is} the system's ontology at any moment---the set of all tokens that have
been observed, with their measurement history.

A new token does not need to be ``added to the vocabulary.'' It is hashed, its bit is
set, and its entry in the LUT is created or incremented. The ontology grows
organically with every measurement event.

% ==========================================================================
\section{Principle 2: IICA Morphisms as Relational Foundation}
\label{sec:iica}
% ==========================================================================

\begin{principle}[IICA Relations]
\label{prin:iica}
Every relation between elements in an EWM is an \textbf{IICA morphism}: Idempotent,
Immutable, and Content-Addressed. Composition of IICA morphisms is IICA. This single
theorem eliminates the need for coordination, consensus, or naming at any scale of
system composition.
\end{principle}

\subsection{The Three Properties}

\begin{definition}[IICA Morphism]
A function $f: X \to Y$ is an IICA morphism iff:
\begin{itemize}[leftmargin=*,nosep]
  \item \textbf{Idempotent:} For every $x \in X$, computing $f(x)$ repeatedly---invoking
    $f$ on the same argument $x$ any number of times---always yields the same output
    $f(x)$.  When the codomain coincides with the domain ($X = Y$), this specializes to
    the classical algebraic form $f(f(x)) = f(x)$ (e.g., HLLSet union:
    $A \cup A = A$).  When $X \neq Y$ (e.g., a hash function
    $h: \text{Tokens} \to \{0,1\}^{160}$), the output type does not match the input type,
    so self-composition $h(h(x))$ is not meaningful; idempotence reduces to
    \textbf{determinism}: the same token $a$ hashes to the same digest $h(a)$ every time,
    regardless of when, where, or how many times $h$ is invoked.
  \item \textbf{Immutable:} $f(x) = y$ is fixed once computed. $y$ never changes.
  \item \textbf{Content-Addressed:} If $a = b$ then $f(a) = f(b)$. The output
    \textit{is} the address---given the content, you can find it again without a
    directory or a naming service.
\end{itemize}
\end{definition}

The canonical IICA morphism is a cryptographic hash function. $\text{SHA-1}(x)$ always
produces the same 160-bit digest for the same $x$ (idempotent as determinism); the digest
never changes (immutable); and the digest serves as the content address in IPFS/IPLD
(content-addressed).

\subsection{The Composition Theorem}

\begin{theorem}[IICA Composition]
\label{thm:composition}
Let $f_1, f_2, \ldots, f_n$ be IICA morphisms. Then their composition
$f_n \circ f_{n-1} \circ \cdots \circ f_1$ is an IICA morphism.
\end{theorem}

\begin{proof}
Each $f_i$ is idempotent, immutable, and content-addressed. Composition preserves each
property:
\begin{itemize}[leftmargin=*,nosep]
  \item \textbf{Idempotency:} For any input $x$, the composition
    $F = f_n \circ \cdots \circ f_1$ is a deterministic function: the same $x$ always
    triggers the same chain of intermediate results $f_1(x), f_2(f_1(x)), \ldots$ and
    produces the same final output $F(x)$.  When $X_1 = Y_n$ (the composition is an
    endomorphism), this additionally yields the algebraic form $F(F(x)) = F(x)$.
  \item \textbf{Immutability:} Each intermediate result $f_i(\ldots)$ is fixed once
    computed; the entire chain of intermediates is immutable.
  \item \textbf{Content-Addressability:} $a = b \implies f_1(a) = f_1(b) \implies
    \cdots \implies F(a) = F(b)$.
\end{itemize}
\end{proof}

This theorem is the engine of compositional emergence. It means that \textbf{any
pipeline of IICA morphisms is itself an IICA morphism}. You can compose tokenizers,
hashers, bridges between domains, materializers, and temporal aggregators without
ever introducing mutable state, naming conflicts, or coordination overhead.

\subsection{The IICA Pipeline}

In HLLSet Algebra, the full measurement pipeline is an IICA composition:

\begin{equation}
  \text{Tokens} \xrightarrow{h_1} \text{Token Hashes (u64)}
  \xrightarrow{h_2} \text{Bit Positions} \xrightarrow{h_3} \text{Bridge HLLSet}
  \xrightarrow{h_4} \text{Materialized Tokens} \xrightarrow{h_5} \cdots
\end{equation}

Each $h_i$ is a MurmurHash3-based IICA morphism. The composition $h_5 \circ h_4 \circ
h_3 \circ h_2 \circ h_1$ is IICA. This means:
\begin{itemize}[leftmargin=*,nosep]
  \item \textbf{No consensus protocols.} Two nodes ingesting the same tokens produce
    identical HLLSets. The lattice converges by construction.
  \item \textbf{No versioning.} HLLSets are immutable. ``Version $n+1$'' is a new
    HLLSet with a new CID. The old HLLSet remains valid and addressable.
  \item \textbf{No naming.} Every HLLSet's CID is its SHA-1 hash. No directory, no
    registry, no naming service.
  \item \textbf{Arbitrary nesting.} A swarm of $N$ robots, each with $M$ sensors,
    produces $N \times M$ IICA pipelines. Their composition at the command post is IICA.
    The theory does not change. The operations are identical. Only the scale differs.
\end{itemize}

\subsection{R-Links: IICA Morphisms as Objects}

A critical structural feature of EWM is that IICA morphisms between elements are
themselves elements. In HLLSet Algebra, the R-link $R_{AB} = A \cap B$ is both:
\begin{itemize}[leftmargin=*,nosep]
  \item A \textbf{morphism} from $A$ to $B$ (it captures the structural overlap that
    relates them).
  \item An \textbf{object} (it is an HLLSet with its own CID, storable in IPFS,
    composable with further intersections).
\end{itemize}

This self-referential property---morphisms as objects---is the structural signature of
a category that can model its own relations. It anticipates the Grothendieck topology
target: in a sheaf topos, the subobject classifier is itself an object of the topos.

% ==========================================================================
\section{Principle 3: The Ashby--Bootstrap Correspondence}
\label{sec:ashby}
% ==========================================================================

\begin{principle}[Ashby--Bootstrap Correspondence]
\label{prin:ashby}
Ashby's Requisite Variety Law implies that connected systems must match in
representational capacity. When cardinalities are mismatched, EWM resolves the
imbalance not by reducing capacity but by \textbf{bootstrap permutation}---multiplying
measurements to expand the effective cardinality of the smaller system until it
matches the larger.
\end{principle}

\subsection{Ashby's Law and the Cardinality Problem}

Ashby's Requisite Variety Law \cite{ashby1952} states that a controller must have at
least as much variety as the system it controls. In EWM terms: when two systems are
connected by IICA morphisms, their representational capacity must be balanced.

Consider the HLLSet architecture. The LUT (Lookup Table) maps tokens to bit positions.
For a vocabulary of 80K Chinese characters, the LUT has 80K entries. The HLLSet bitmask
has $1024 \times 32 = 32{,}768$ bit positions. A naive reading of Ashby's law would
suggest that the LUT and HLLSet are mismatched: 80K vs.\ 32K.

But this reading is superficial. The HLLSet does not store individual tokens---it
stores \textit{combinations} of tokens through hash collisions. The bitmask encodes the
OR of all tokens that hash to each position. The information capacity is not determined
by the number of bit positions but by the number of \textit{distinguishable bitmask
configurations}, which is $2^{32,768}$.

\subsection{Bootstrap Permutation}

The deeper problem is not total capacity but \textbf{measurement resolution}. When a
single token maps to a single bit position, the measurement is coarse: the bit is either
set or not, and the materializer cannot distinguish which of several tokens set it.

The solution is \textbf{bootstrap permutation}: multiply the measurements by applying a
family of permutations to the token stream before measurement.

\begin{definition}[Bootstrap Permutation]
Let $\Pi = \{\pi_1, \ldots, \pi_k\}$ be a family of permutations on token sequences.
Given a token stream $S = (t_1, \ldots, t_n)$, the bootstrapped measurement is:
\begin{equation}
  H_{\Pi}(S) = \bigoplus_{j=1}^{k} H(\pi_j(S))
\end{equation}
where $H(\pi_j(S))$ is the HLLSet of the permuted stream and $\oplus$ is bitwise OR.
\end{definition}

Each permutation $\pi_j$ produces a different ``view'' of the same token stream. The
union of these views creates a richer measurement---multiple bit positions per original
token, finer discrimination, and better materialization fidelity.

\subsection{n-Grams as Bootstrap}

In the current HLLSet implementation, \textbf{n-grams} are the concrete bootstrap
family. Given a token stream $(t_1, t_2, t_3, t_4, \ldots)$, the 2-gram view is
$((t_1, t_2), (t_2, t_3), (t_3, t_4), \ldots)$ and the 3-gram view is
$((t_1, t_2, t_3), (t_2, t_3, t_4), \ldots)$. Each n-gram is hashed as a compound
token, producing a distinct set of bit positions.

\begin{proposition}[n-Gram Cardinality Expansion]
A vocabulary of $|V|$ tokens, measured with n-grams up to order $K$, produces up to
$\sum_{n=1}^{K} |V|^n$ distinct measurement points. For $|V| = 80{,}000$ and $K = 3$,
this is approximately $5.12 \times 10^{14}$ possible n-gram tokens, massively exceeding
the $32{,}768$ HLLSet bit positions.
\end{proposition}

The bootstrap permutation bridges the cardinality gap: a vocabulary that appears
``too small'' relative to the HLLSet is expanded through n-gram composition; a
vocabulary that appears ``too large'' is aggregated through hash collisions into the
fixed bit space. The balance is achieved through \textit{measurement multiplication},
not cardinality reduction.

\subsection{General Bootstrap Families}

n-grams are one family of bootstrap permutations. The principle generalizes:
\begin{itemize}[leftmargin=*,nosep]
  \item \textbf{Skip-grams:} $(t_i, t_{i+k})$ for various skip distances $k$.
  \item \textbf{Convolutional windows:} weighted combinations of neighboring tokens.
  \item \textbf{Spectral permutations:} frequency-domain rearrangements via FFT.
  \item \textbf{Random projections:} a family of random hash functions, each producing
    a distinct HLLSet view.
  \item \textbf{Cross-modal bootstraps:} the same stream measured through different
    sensory modalities, each with its own IICA pipeline.
\end{itemize}

The bootstrap family is a design parameter of the EWM. Different families produce
different emergent ontologies from the same underlying token stream. The choice of
family is the EWM analog of architecture selection in neural networks---it determines
what patterns the system can recognize.

% ==========================================================================
\section{Principle 4: Noether's Theorem as EWM Evolution Law}
\label{sec:noether}
% ==========================================================================

\begin{principle}[Noether Evolution]
\label{prin:noether}
In an EWM, every continuous symmetry of the measurement regime yields a conserved
quantity. Structural change is governed by conservation laws, not by gradient descent.
The system evolves by redistributing information across the lattice while preserving
global invariants.
\end{principle}

\subsection{From Physics to World Models}

In physics, Noether's theorem \cite{noether1918} establishes that every differentiable
symmetry of the action of a physical system corresponds to a conservation law:
time-translation symmetry $\rightarrow$ energy conservation; spatial-translation
symmetry $\rightarrow$ momentum conservation; rotational symmetry $\rightarrow$ angular
momentum conservation.

In an EWM, the ``action'' is the measurement regime---the accumulation of tokens into
HLLSets across temporal layers. The symmetries are transformations of the measurement
process that leave the structural relations invariant. The conserved quantities are
global properties of the lattice that persist across evolution steps.

\subsection{The Temporal Pyramid and [UM]-Net Cycles}

Before stating the invariant, we briefly introduce the two temporal structures that give
it physical meaning.

\textbf{The temporal pyramid} is a configurable $K$-layer compression stack (default
$K = 7$: second, minute, hour, day, week, month, year). Each layer $L_i$ accumulates
all measurement events within its temporal window via monotonic union:
$L_i = \bigcup S(t)$ for $t$ in the window. At each window boundary, the layer cascades
into the next coarser layer: $L_{i+1} \leftarrow L_{i+1} \cup L_i$, then $L_i$ resets.
After compression, layers are mutually exclusive---no time slice appears in more than one
layer. The complete system state is their union:
$H_{\text{system}} = \bigcup_{i=0}^{K} L_i$. Compression ratios range from $60\!:\!1$
(seconds$\to$minutes) to $12\!:\!1$ (months$\to$years); seven layers compress
approximately $3.15 \times 10^7$ seconds into one year-layer HLLSet.

\textbf{[UM] agents and [UM]-Net} provide a complementary temporal structure at shorter
timescales.  A \textbf{[UM]} (unified model) is the concrete implementation of a single
EWM: a processing node that receives a stream of materialized encodings from the external
environment \texttt{[E]}, ingests them into its internal HLLSet lattice, and materializes
the result back to a disambiguated encoding stream for downstream consumption.  Its
internal pipeline is a compound IICA morphism:

\begin{equation}
  \text{[UM]} \;=\;
  \text{(encodings)} \xrightarrow{\text{ingest}} \{\text{HLLSets}\}
  \xrightarrow{\text{materialize}} \text{(encodings)}
\end{equation}

The ingestion step hashes the incoming encodings into bit-vectors and accumulates them
into the lattice via monotonic union; the materialization step projects the lattice
through the agent's LUT and recovers ordered, disambiguated tokens via De Bruijn
reconstruction (Section~\ref{sec:hllset}).  Each individual step in the pipeline
(hashing, union, LUT projection) is IICA.  The [UM] node as a whole is IICA
\textit{between commits}: as long as its internal TF-Vec and HLLSet lattice are
fixed, the same input encodings always produce the same output encodings.  State
changes---new measurements accumulated, ranks updated---occur only at explicit
\textbf{commit events}, which define transactional boundaries.  Within a transaction,
[UM] behaves as a pure IICA morphism; across transactions, the same input may produce
different output as the lattice evolves.

A \textbf{[UM]-Net} is a directed graph $G = (A, E)$ of [UM] nodes connected by
fire-and-forget edges.  The external environment \texttt{[E]} feeds encodings into source
nodes; each [UM] processes independently and fires its output downstream with no
acknowledgment, no handshake, and no retry.  Convergence is guaranteed by IICA and CRDT
union at confluence nodes.  The full system diagram is:

\begin{equation}
  \texttt{[E]} \longrightarrow \text{(encodings)} \longrightarrow
  \texttt{[UM]}_1 \longrightarrow \texttt{[UM]}_2 \longrightarrow \cdots
\end{equation}

Critically, \textbf{cycles are not prohibited}: a recurrent loop
$a \to \cdots \to a$ in $G$ is a \textit{short-term memory} structure, resolved by
temporal separation---$\text{out}_a(t)$ feeds back as $\text{in}_a(t + k)$ after $k$
ingestion steps around the cycle.  Each pass through the cycle produces a new observation
at a later time, not a re-entry of the same data.  The evolution equation
$H(t) = H(S(t), H(t-1), D, R, N)$ captures both pyramid cascade and cycle recurrence
under a single dynamics.

\subsection{The Union Invariant}

The fundamental Noether symmetry in EWM is \textbf{path independence}: information can
travel through the temporal pyramid via multiple routes (L0$\to$L1$\to$L2, or
L0$\to$L2 directly, or L1$\to$L4$\to$L5). The conserved quantity is the union of all
temporal layers:

\begin{equation}
  H_{\text{system}}(t) = \bigcup_{i=0}^{K} L_i(t) = \text{constant over time}
\end{equation}

\begin{theorem}[Temporal Union Invariant]
For any sequence of measurement events and any cascade policy, the bitwise union of all
temporal layers is monotonic and convergent. Information is never lost from
$H_{\text{system}}$; it is only redistributed across layers.
\end{theorem}

This is eventual consistency by construction---not as a protocol to implement, but as a
mathematical property of the lattice. No consensus algorithm. No leader election. No
retry logic. The lattice converges because the union is monotonic and multiple paths
guarantee that every bit reaches every layer that needs it.

\subsection{The DRN Conservation Law}

The DRN decomposition (Departed, Retained, Novel) expresses evolution as a conservation
balance:

\begin{equation}
  H(t) = H(S(t), H(t-1), D(t-1), R(t-1), N(t))
\end{equation}

where:
\begin{itemize}[leftmargin=*,nosep]
  \item $S(t)$ --- current observation as HLLSet
  \item $H(t-1)$ --- previous system state
  \item $D(t-1) = H(t-1) \setminus H(t)$ --- departed context
  \item $R(t-1) = H(t-1) \cap H(t)$ --- retained context
  \item $N(t) = H(t) \setminus H(t-1)$ --- novel context
\end{itemize}

The Noether steering condition is a conservation law on bit count across the DRN
transition:

\begin{equation}
  |\,\text{card}(N(t)) - \text{card}(D(t-1))\,| \to 0
\end{equation}

At steady state, the number of new bits entering the system equals the number of old
bits departing. The structural identity of the system is preserved even as its content
evolves. In rank-weighted form:

\begin{equation}
  \left|\sum_{b \in N(t)} R_b(t) - \sum_{b \in D(t-1)} R_b(t-1)\right| \to 0
\end{equation}

This is the EWM analog of energy conservation: the total ``rank mass'' of the system is
conserved across transitions. Learning is not accumulation---it is redistribution.

\subsection{The Holographic Property as Noether Consequence}

The holographic property---that the lattice top $H_{\text{system}}$ plus the TF stack
can reconstruct any past state---is a direct consequence of the union invariant:

\begin{equation}
  \text{past\_state}(t) \approx H_{\text{system}}(\text{now}) \odot
  \text{TF}_{\text{stack}}[t]
\end{equation}

Because $H_{\text{system}}$ never loses information (union invariant), and the TF stack
records \textit{when} each bit position was active (temporal distribution), the
reconstruction is always possible. The information is not stored separately for each
time step---it is encoded holographically in the lattice, with the TF stack serving as
the reference beam for temporal selectivity.

\subsection{The Fisher Matrix as Structural Derivative}

The cross-layer Fisher-like matrix captures the second-order structure of evolution:

\begin{equation}
  F_{bb'} = \sum_{i=0}^{K} B^{(i)}_b \cdot B^{(i)}_{b'}
\end{equation}

where $B^{(i)}_b = 1$ if bit $b$ is set in layer $i$. $F_{bb'}$ counts how many layers
have both bits $b$ and $b'$ set simultaneously. High $F_{bb'}$ indicates systemic
co-occurrence (a structural invariant). Low $F_{bb'}$ indicates noise. The Fisher matrix
enables the Noether controller to distinguish isolated fluctuations from systemic phase
transitions.

\subsection{Noether Controller}

The Noether controller monitors the cross-layer relationship matrix and decides which
temporal layer should drive the next action. When the BSS between L0 (current second)
and L1 (recent minute) drops below a threshold, the controller detects divergence and
overrides instinct with broader context. When L0 diverges from L3 (day-scale goals),
the controller triggers re-routing.

The controller is the \textbf{structural actuator} of the Noether principle. It does not
adjust weights---it selects which layer of accumulated measurement history is relevant
to the current situation. The six control parameters (pyramid depth, rank functions,
attention threshold, steering thresholds) are the knobs of structural evolution, not the
parameters of a learned function.

% ==========================================================================
\section{HLLSet Algebra: Concrete Instantiation}
\label{sec:hllset}
% ==========================================================================

The four principles above are not abstract. They are realized in the HLLSet Algebra
framework \cite{astesj2026}, whose components map directly onto the EWM structure.

\subsection{The Patch Space as Measurement Site}

The HLLSet bitmask is a \textbf{holographic plate} of $32{,}768$ patches (bit
positions), each recording a binary mark: 0 = untouched, 1 = at least one token hashed
here. The hash function $h: \mathcal{T} \to \{0,\ldots,32767\}$ is the IICA measurement
morphism. The LUT $\mathcal{L}$ tracks the inverse image of $h$, enabling materialization
(recovering which tokens could have set which bits).

This is the minimal EWM: a finite set of measurement sites, an IICA morphism from tokens
to patches, and a monotonic accumulation operation (bitwise OR) that makes each patch a
join-semilattice. Principle~1 (emergent ontology) is realized by the LUT. Principle~2
(IICA relations) is realized by the hash pipeline and R-link lattice operations.

\subsection{The Temporal Pyramid as Noether Structure}

The 7-layer temporal pyramid (second, minute, hour, day, week, month, year) implements
Principle~4 (Noether evolution). Each layer accumulates measurements over its temporal
window via union. The cascade policy (L0 overflows into L1, L1 into L2, etc.) is a
discrete-time dynamical system whose conserved quantity is the union invariant.

\subsection{n-Grams as Bootstrap}

The tokenizer's n-gram pipeline (1-grams, 2-grams, 3-grams with boundary padding)
implements Principle~3 (bootstrap permutation). Each n-gram order is a distinct
measurement view of the same token stream. The union of n-gram HLLSets produces a richer
representation than any single order alone.

\subsection{Rank Algebra as Learning}

Learning is rank rearrangement. HLLSets are immutable (IICA). The Forth dictionary
assigns ranks to HLLSets based on the current TF vector. A rank adjustment is a
structural change (which HLLSet drives action) without any parametric change (the
HLLSets themselves are untouched). This satisfies Principle~4: learning is
redistribution of attention across the lattice, not accumulation of weight updates.

% ==========================================================================
\section{Toward Grothendieck Topology}
\label{sec:gt}
% ==========================================================================

The four principles above point toward a deeper mathematical structure:
\textbf{Grothendieck topology} \cite{grothendieck1972} on the category of measurement
patches. We articulate this target not as a completed formalism, but as a compass
bearing---the ``spirit'' in which the EWM framework is developed.

\subsection{The Site: A Category of Patches}

Let $\mathcal{P}$ be a finite discrete category whose objects are patch addresses and
whose only morphisms are identities. This is the \textbf{site}---the space of measurement
locations.

A \textbf{presheaf} on $\mathcal{P}$ assigns to each patch $p$ a set of possible marks
$\Sigma(p)$ (the mark alphabet) and to each morphism (only identities) the trivial
restriction. An HLLSet is a global section of this presheaf:

\begin{equation}
  H \in \Gamma(\mathcal{P}, \Sigma) = \prod_{p \in \mathcal{P}} \Sigma(p)
\end{equation}

For standard HLLSets, $\Sigma(p) = \{0, 1\}$ (a bit) with accumulation $\oplus =
\text{OR}$. For the generalized case (GENERAL\_HLLSET\_ALGEBRA.md), $\Sigma(p)$ can be
any finite join-semilattice: $\{0, 1, 2\}$ (trits with max), $\{0, 1, 2, 3\}$ (quits),
or arbitrary discrete alphabets.

\subsection{Coverages and the Gluing Condition}

A Grothendieck topology on $\mathcal{P}$ specifies which families of morphisms
$\{U_i \to U\}$ count as \textbf{covering families}. The gluing condition states: if
data on each $U_i$ is consistent on overlaps $U_i \times_U U_j$, then there exists
unique data on $U$ that restricts to the given data on each $U_i$.

In EWM terms, the covering families are the \textbf{bootstrap permutation family}
(Principle~3). Each permutation $\pi_j$ provides a partial view of the token stream
(a measurement on a ``sub-site''). The gluing condition is satisfied when these partial
views are consistent---when the same token maps to compatible marks across different
permutations.

The HLLSet union $\bigcup_{\pi \in \Pi} H_\pi$ is the \textbf{sheafification} of the
bootstrap family: it combines all partial measurements into a global section.

\subsection{Geometric Morphisms as Universal Bridges}

A geometric morphism $f: \text{Sh}(\mathcal{P}) \to \text{Sh}(\mathcal{Q})$ between
sheaf toposes on different patch spaces is the categorical analog of the Universal
Bridge. It consists of:
\begin{itemize}[leftmargin=*,nosep]
  \item A \textbf{direct image} $f_*$ that pushes HLLSets from $\mathcal{P}$ to
    $\mathcal{Q}$.
  \item An \textbf{inverse image} $f^*$ that pulls HLLSets from $\mathcal{Q}$ back to
    $\mathcal{P}$.
\end{itemize}

In the current implementation, the Universal Bridge implements the direct image via
re-representation: format each bit position as a token, re-hash into the target space.
The inverse image (materialization from target back to source) is approximate---it
depends on the LUT and is subject to the Statistics Constraint (TF vectors are not
transferred across bridges).

\subsection{The Subobject Classifier and R-Links}

In a topos, the \textbf{subobject classifier} $\Omega$ classifies subobjects: for any
object $X$, there is a bijection between subobjects of $X$ and morphisms $X \to \Omega$.
For the Boolean topos on $\mathcal{P}$, $\Omega$ is the two-element set $\{0, 1\}$
with the OR accumulation. Every HLLSet $H$ is a morphism $H: \mathcal{P} \to \Omega$.

The R-link $R_{AB} = A \cap B$ is the \textbf{pullback} of $A$ and $B$ over $\Omega$.
In categorical terms:
\begin{equation}
  \begin{CD}
    R_{AB} @>>> A \\
    @VVV @VVV \\
    B @>>> \Omega
  \end{CD}
\end{equation}

This is why R-links are both morphisms and objects: they are pullbacks in the topos.
The lattice of HLLSets is a \textbf{Heyting algebra}---the internal logic of the topos.

\subsection{What the Grothendieck Target Gives Us}

Adopting Grothendieck topology as the target formalism would provide:

\begin{enumerate}[leftmargin=*,nosep]
  \item \textbf{A rigorous gluing theorem.} Given consistent measurements on a covering
    family, the global state is uniquely determined. This formalizes Principle~3 (why
    bootstrap permutations are sufficient).
  \item \textbf{A systematic theory of bridges.} Geometric morphisms between patch
    topoi classify all possible Universal Bridges. Different choices of direct/inverse
    image correspond to different bridge implementations.
  \item \textbf{An internal logic.} The Heyting algebra structure on the HLLSet lattice
    is the internal logic of the topos. Union, intersection, and implication correspond
    to logical OR, AND, and IMPLIES. The EWM \textit{reasons} within its own lattice.
  \item \textbf{Cohomology.} Sheaf cohomology groups would measure the ``obstruction''
    to gluing---quantifying how much information is lost when different bootstrap views
    disagree.
\end{enumerate}

We are not yet ready to formalize EWM as a full Grothendieck topology. The current
framework is a \textbf{pre-sheaf structure} on a discrete site with a bootstrap covering
family. The four principles---emergent ontology, IICA morphisms, bootstrap permutation,
and Noether evolution---are the structural requirements that any such topology must
satisfy. The destination may differ from the current compass bearing. But the bearing
itself---gluing local measurements into global states through consistent covering
families---is the ``spirit'' that guides the framework.

% ==========================================================================
\section{Discussion}
\label{sec:discussion}
% ==========================================================================

\subsection{Unification}

The four principles unify the existing body of work:

\begin{itemize}[leftmargin=*,nosep]
  \item \textbf{HLLSet Theory} \cite{astesj2026} establishes the base lattice
    (Principle~1, 2) and the IICA gate (Principle~2).
  \item \textbf{The Self-Reprogramming Architecture} and \textbf{STANDARD.md} specify
    the temporal pyramid, DRN decomposition, Noether controller, and holographic memory
    (Principle~4).
  \item \textbf{GENERAL\_HLLSET\_ALGEBRA.md} identifies the categorical structure
    (Bool$_{32768}$, IICA\_Hash, Patch category) and the Grothendieck target
    (Section~\ref{sec:gt}).
  \item \textbf{The n-gram tokenizer} and \textbf{Universal Bridge} implement
    bootstrap permutation (Principle~3).
  \item \textbf{CAAL-LLM} validates that the principles produce working intelligence:
    80\% accuracy on driving scenarios from 10 sentences, no gradient descent
    \cite{astesj2026}.
\end{itemize}

\subsection{What EWM Is Not}

It is important to state what EWM is \textit{not}:

\begin{itemize}[leftmargin=*,nosep]
  \item \textbf{Not a neural network.} There are no weight matrices, no activation
    functions, no backpropagation. The lattice is the computation; ranks are the
    attention; the dictionary is the memory.
  \item \textbf{Not a token-based model.} There is no fixed vocabulary. Elements emerge
    from measurement accumulation. A new token is hashed and enters the lattice without
    architectural change.
  \item \textbf{Not a probabilistic model in the Bayesian sense.} There are no
    priors, no posteriors, no sampling. The BSS (Bell State Similarity) is a
    deterministic bitwise measure, not a probability.
  \item \textbf{Not a database.} Although content-addressed and IPFS-native, the
    lattice is a computational structure, not a storage system. The HLPP protocol
    specifies persistence, but the lattice operations are the computation.
\end{itemize}

\subsection{Open Questions}

\begin{enumerate}[leftmargin=*,nosep]
  \item \textbf{Optimal bootstrap families.} What is the minimal set of permutations
    that guarantees $\epsilon$-accurate materialization for a given token stream
    distribution? The current n-gram choice is empirical.
  \item \textbf{Convergence rates.} Under what conditions does the Noether controller
    converge to a stable layer selection policy? What is the relationship between
    environmental volatility and convergence time?
  \item \textbf{Topos formalization.} What is the minimal Grothendieck topology that
    captures the bootstrap covering condition? Is the current discrete site sufficient,
    or do we need a non-trivial coverage?
  \item \textbf{Cross-modal emergence.} If tokens come from different sensory modalities
    (text, image patches, audio spectra), how does the unified lattice represent
    cross-modal elements? The Universal Bridge provides translation; emergence across
    modalities is unexplored.
  \item \textbf{The Statistics Constraint.} Why must TF vectors NOT be transferred
    across bridges? The constraint is empirical (LUT initialization failure). A
    categorical explanation---TF as a sheaf of local statistics that does not push
    forward under geometric morphisms---is plausible but unproven.
\end{enumerate}

\subsection{The Road Ahead}

The EWM framework is at the beginning of its development. The four principles provide
a structural compass. The HLLSet Algebra implementation provides a concrete testbed.
The Grothendieck target provides theoretical direction. The path forward involves:

\begin{enumerate}[leftmargin=*,nosep]
  \item \textbf{Harden the implementation.} Complete the FPGA backend, distributed mesh
    networking, and self-ingestion pipeline to validate the principles at scale.
  \item \textbf{Explore the bootstrap design space.} Systematic experiments with
    different permutation families, mark alphabets (trits, quits), and accumulation
    semilattices.
  \item \textbf{Develop the categorical theory.} Move from pre-sheaf to sheaf, from
    discrete site to Grothendieck topology, from Bridge to geometric morphism.
  \item \textbf{Bridge to applications.} The CAAL-LLM proof-of-concept demonstrates
    viability. Real-time adaptive systems (robotics, autonomous vehicles, industrial
    control) are the natural deployment targets for an FPGA-native, gradient-free world
    model.
\end{enumerate}

% ==========================================================================
% ── References ────────────────────────────────────────────────────────────
% ==========================================================================

\begin{thebibliography}{9}

\bibitem{astesj2026}
A.~Mylnikov.
\emph{HLLSet Theory: A Unified Framework for Probabilistic Knowledge Representation}.
ASTESJ, Vol.~11, No.~2, pp.~62--74, 2026.
\url{https://www.astesj.com/publications/ASTESJ_110202.pdf}

\bibitem{sgs2024}
A.~Mylnikov.
\emph{Self Generative Systems (SGS) and Its Integration with AI Models}.
AISNS '24: Proceedings of the 2024 2nd International Conference on Artificial
Intelligence, Systems and Network Security, pp.~345--354, 2024.
\url{https://doi.org/10.1145/3714334.3714392}

\bibitem{ashby1952}
W.~R.~Ashby.
\emph{Design for a Brain}.
Chapman \& Hall, 1952.

\bibitem{noether1918}
E.~Noether.
\emph{Invariante Variationsprobleme}.
Nachr.~d.~K\"onig.~Gesellsch.~d.~Wiss.~zu G\"ottingen, Math-phys.~Klasse,
pp.~235--257, 1918.

\bibitem{flajolet2007}
P.~Flajolet, \'{E}.~Fusy, O.~Gandouet, F.~Meunier.
\emph{HyperLogLog: the analysis of a near-optimal cardinality estimation algorithm}.
AofA, 2007.

\bibitem{lecun2022}
Y.~LeCun.
\emph{A Path Towards Autonomous Machine Intelligence}.
Open Review, 2022.
\url{https://openreview.net/forum?id=BZ5a1r-kVsf}

\bibitem{grothendieck1972}
A.~Grothendieck, J.~L.~Verdier.
\emph{Th\'eorie des Topos et Cohomologie \'Etale des Sch\'emas} (SGA 4).
Lecture Notes in Mathematics, Vols.~269, 270, 305. Springer, 1972--1973.

\bibitem{maclane1992}
S.~Mac Lane, I.~Moerdijk.
\emph{Sheaves in Geometry and Logic: A First Introduction to Topos Theory}.
Springer, 1992.

\end{thebibliography}

% == Acknowledgment =========================================================
\section*{Acknowledgment}
This consolidation was developed in dialogue with DeepSeek Code, an AI research
assistant. The four principles---Emergent Ontology, IICA Morphisms, Ashby--Bootstrap
Correspondence, and Noether Evolution---crystallized through structured dialogue between
Alex Mylnikov and DeepSeek during June--August 2026. The core HLLSet Algebra, IICA
principles, and the ASTESJ theoretical foundation were established prior to this
collaboration. The Grothendieck topology target was identified in
\texttt{GENERAL\_HLLSET\_ALGEBRA.md} (July--August 2026). This document is a
consolidation of research documented in \texttt{\_DOCS/dev/} and
\texttt{\_DOCS/proposals/} of the \texttt{hllset-next} repository.

\end{document}
